\section{Introduction to Convection Heat Transfer}

\subsection{What is Convection Heat Transfer?}
Convection is the mode of energy transfer between a solid surface and the adjacent liquid or gas that is in motion. It involves two mechanisms:
\begin{enumerate}
    \item \textbf{Energy transfer due to random molecular motion (conduction):} Occurs within the fluid.
    \item \textbf{Energy transfer by bulk motion (advection):} The fluid carries energy with it as it moves.
\end{enumerate}
Convection heat transfer is highly dependent on the fluid properties (dynamic viscosity $\mu$, thermal conductivity $\kth$, density $\rho$, and specific heat $\Cp$), as well as fluid velocity. It also depends on the geometry and roughness of the solid surface, and the type of fluid flow (e.g., laminar or turbulent).

\subsection{Types of Convection}
\begin{itemize}
    \item \textbf{Forced Convection:} The fluid motion is induced by external means, such as a fan, pump, or the wind.
    \item \textbf{Natural (or Free) Convection:} The fluid motion is caused by buoyancy forces, which arise from density differences in the fluid due to temperature variations. Warmer, less dense fluid rises, and cooler, denser fluid falls.
\end{itemize}

\subsection{The No-Slip Condition}
When a fluid flows over a solid surface, the layer of fluid directly in contact with the surface "sticks" to it. This means the fluid velocity at the surface is zero relative to the surface. This is known as the no-slip condition. Consequently, heat transfer at the solid-fluid interface is purely by conduction. The overall heat transfer, however, is convective due to the subsequent mixing and bulk motion of the fluid.

\section{Boundary Layers}
When fluid flows over a surface, various boundary layers develop where the fluid properties are significantly affected by the presence of the surface.

\subsection{Velocity Boundary Layer}
\begin{itemize}
    \item \textbf{Definition:} The velocity boundary layer is the region of fluid flow over a surface where the velocity of the fluid changes from zero at the surface (due to the no-slip condition) to the free-stream velocity (undisturbed flow velocity).
    \item \textbf{Regions:}
        \begin{itemize}
            \item \textbf{Laminar Boundary Layer:} Fluid flow is smooth and orderly.
            \item \textbf{Transition Region:} Flow transitions from laminar to turbulent.
            \item \textbf{Turbulent Boundary Layer:} Fluid flow is chaotic and characterized by eddies and random fluctuations. Near the surface, a very thin "laminar sublayer" or "viscous sublayer" still exists.
        \end{itemize}
    \item \textbf{Thickness ($\delta$):} The velocity boundary layer thickness is typically defined as the distance from the surface where the fluid velocity reaches 99\% of the free-stream velocity. It grows in the direction of flow.
\end{itemize}

\subsection{Thermal Boundary Layer}
\begin{itemize}
    \item \textbf{Definition:} The thermal boundary layer is the region of fluid flow over a surface where the temperature of the fluid changes from the surface temperature to the free-stream temperature.
    \item \textbf{Thickness ($\delta_T$):} The thermal boundary layer thickness is defined as the distance from the surface where the fluid temperature reaches 99\% of the free-stream temperature difference, i.e., $\T - \T_s = 0.99(\Tinf - \T_s)$. It also grows in the direction of flow.
    \item \textbf{Relationship to Velocity Boundary Layer:} The relative thicknesses of the velocity and thermal boundary layers are governed by the Prandtl number.
\end{itemize}

\section{Dimensionless Numbers in Convection}
Dimensionless numbers are crucial in heat transfer as they allow for generalization of experimental data and provide insight into the relative importance of different physical phenomena.

\subsection{Reynolds Number ($\ReNum$)}
\begin{itemize}
    \item \textbf{Definition:} A dimensionless number that quantifies the ratio of inertial forces to viscous forces in a fluid flow.
    \item \textbf{Significance:} It is used to predict flow patterns in different fluid flow situations and determine whether the flow is laminar or turbulent.
        \begin{itemize}
            \item Low $\ReNum$ (e.g., $< 2300$ for pipe flow): Laminar flow
            \item High $\ReNum$ (e.g., $> 4000$ for pipe flow): Turbulent flow
        \end{itemize}
\end{itemize}

\subsection{Prandtl Number ($\PrNum$)}
\begin{itemize}
    \item \textbf{Definition:} A dimensionless number that expresses the ratio of momentum diffusivity (kinematic viscosity) to thermal diffusivity.
    \item \textbf{Significance:} It characterizes the relative thickness of the velocity boundary layer and the thermal boundary layer.
        \begin{itemize}
            \item $\PrNum \approx 1$: Momentum and thermal diffusivities are similar, so boundary layers have similar thicknesses ($\delta \approx \delta_T$).
            \item $\PrNum > 1$: Momentum diffusivity dominates, so velocity boundary layer is thicker than thermal boundary layer ($\delta > \delta_T$). (Typical for oils and liquids)
            \item $\PrNum < 1$: Thermal diffusivity dominates, so thermal boundary layer is thicker than velocity boundary layer ($\delta < \delta_T$). (Typical for liquid metals)
        \end{itemize}
\end{itemize}

\subsection{Nusselt Number ($\NuNum$)}
\begin{itemize}
    \item \textbf{Definition:} A dimensionless number that represents the ratio of convective heat transfer to conductive heat transfer across the fluid layer.
    \item \textbf{Significance:} It is a measure of the enhancement of heat transfer from a surface due to convection relative to pure conduction. A larger Nusselt number indicates more effective convection.
\end{itemize}

\subsection{Stanton Number ($\StNum$)}
\begin{itemize}
    \item \textbf{Definition:} A dimensionless number that represents the ratio of heat transfer rate to the heat capacity rate of the flow. It is often used in situations where there is simultaneous heat and momentum transfer.
    \item \textbf{Relationship:} It is related to the Nusselt and Reynolds numbers.
\end{itemize}

\section{Classification of Fluid Flow Types}
Fluid flow can be classified based on various characteristics:
\begin{itemize}
    \item \textbf{Viscous vs. Inviscid:} Whether viscous effects are significant (viscous) or negligible (inviscid).
    \item \textbf{Internal vs. External Flow:} Flow confined by solid surfaces (internal, e.g., pipes) vs. flow over surfaces in an unbounded medium (external, e.g., flow over a flat plate).
    \item \textbf{Compressible vs. Incompressible Flow:} Whether the fluid density changes significantly during flow. For liquids, flow is usually incompressible. For gases, it depends on velocity (incompressible for low Mach numbers).
    \item \textbf{Laminar vs. Turbulent Flow:} Characterized by orderly, smooth streamlines (laminar) or chaotic, fluctuating motion (turbulent).
    \item \textbf{Natural (Free) vs. Forced Convection:} As defined above.
    \item \textbf{Steady vs. Unsteady (Transient) Flow:} Whether flow characteristics (velocity, temperature) change with time.
    \item \textbf{One, Two, or Three-Dimensional Flow:} Number of spatial coordinates needed to describe the flow.
\end{itemize}

\section{Governing Equations for Fluid Flow and Heat Transfer}
For complex fluid flow and heat transfer problems, the fundamental conservation laws are expressed as partial differential equations.

\subsection{Continuity Equation (Conservation of Mass)}
States that mass is conserved in a flow system. For incompressible flow, it implies conservation of volume.

\subsection{Momentum Equations (Navier-Stokes Equations - Conservation of Momentum)}
These equations describe the motion of viscous fluids and are based on Newton's second law. They express the balance between inertial, pressure, viscous, and body forces acting on a fluid element.

\subsection{Energy Equation (Conservation of Energy)}
This equation describes the transport of thermal energy in the fluid, accounting for conduction, convection, and energy generation or dissipation within the fluid.

\section{Lumped System Analysis (Transient Convection)}
\subsection{Concept and Assumptions}
Lumped system analysis is a simplified approach for transient heat transfer that assumes the temperature within a solid body is spatially uniform at any given time. This assumption is valid when the internal temperature gradients within the body are negligible compared to the temperature difference between the body and its surroundings.

\subsection{Characteristic Length and Biot Number}
The validity of the lumped system assumption is determined by the Biot number.
\begin{itemize}
    \item \textbf{Characteristic Length ($\Lc$):} A parameter used to nondimensionalize heat transfer problems. For lumped system analysis, it is defined as the ratio of the body's volume to its surface area: $\Lc = V/\A_s$.
    \item \textbf{Biot Number ($\Bi$):} A dimensionless number that represents the ratio of the conduction resistance within the body to the convection resistance at the body's surface.
        \begin{itemize}
            \item $\Bi = \frac{\hC \Lc}{\kth_{\text{solid}}}$
            \item If $\Bi \ll 0.1$, the lumped system analysis is generally applicable, meaning the temperature within the solid is nearly uniform. This implies that conduction within the solid is much faster than convection at its surface.
            \item If $\Bi \ge 0.1$, internal temperature gradients are significant, and a more detailed transient conduction analysis is required.
        \end{itemize}
\end{itemize}

\subsection{Governing Equation and Solution}
The energy balance for a lumped system states that the rate of convection heat transfer into the body equals the rate of increase of internal energy of the body. This leads to a first-order ordinary differential equation whose solution gives the temperature of the body as a function of time.

\section{Heat Transfer Applications (Convection Examples)}
Convection principles are critical in:
\begin{itemize}
    \item \textbf{Heat Exchangers:} Design and performance of devices that transfer heat between two or more fluids.
    \item \textbf{Electronic Cooling:} Dissipation of heat from computer chips and other electronic components.
    \item \textbf{HVAC Systems:} Heating, ventilation, and air conditioning in buildings.
    \item \textbf{Boiling and Condensation:} Phase-change heat transfer processes involving convection.
    \item \textbf{Atmospheric and Oceanic Flows:} Weather patterns, ocean currents.
\item \textbf{Human Body Thermoregulation:} Heat loss/gain through convection with ambient air.
\end{itemize}

 % Start formulas on a new page for clarity
\section{Formulas Summary for Convection}

\subsection{Fundamental Convection Relations}
\subsubsection*{Newton's Law of Cooling}
\begin{equation}
    \Qdot_{\text{conv}} = \hC \A_s (\T_s - \Tinf)
\end{equation}
where: \\
$\Qdot_{\text{conv}}$ = Convection heat transfer rate (W) \\
$\hC$ = Convection heat transfer coefficient (W/m$^2$$\cdot$K) \\
$\A_s$ = Surface area for convection (m$^2$) \\
$\T_s$ = Surface temperature (K or $^\circ$C) \\
$\Tinf$ = Ambient fluid temperature (K or $^\circ$C)

\subsubsection*{Surface Shear Stress}
\begin{equation}
    \tauS = \muF \frac{\partial u}{\partial y}\Big|_{y=0}
\end{equation}
where: \\
$\tauS$ = Surface shear stress (N/m$^2$) \\
$\muF$ = Dynamic viscosity of the fluid (Pa$\cdot$s or kg/m$\cdot$s) \\
$u$ = Local fluid velocity in the x-direction (m/s) \\
$y$ = Coordinate perpendicular to the surface (m)

\subsubsection*{Drag Force (on a flat plate)}
\begin{equation}
    F_D = \Cf \A_s \frac{\rhoD \Vavg^2}{2}
\end{equation}
where: \\
$F_D$ = Drag force (N) \\
$\Cf$ = Average friction coefficient (dimensionless) \\
$\A_s$ = Surface area of the plate (m$^2$) \\
$\rhoD$ = Fluid density (kg/m$^3$) \\
$\Vavg$ = Free-stream velocity of the fluid (m/s)

\subsection{Dimensionless Numbers}
\subsubsection*{Reynolds Number ($\ReNum$)}
\begin{equation}
    \ReNum_L = \frac{\rhoD \Vavg L}{\muF} = \frac{\Vavg L}{\nuF}
\end{equation}
where: \\
$\ReNum_L$ = Reynolds number (dimensionless, based on characteristic length $L$) \\
$\Vavg$ = Characteristic fluid velocity (m/s) \\
$L$ = Characteristic length (e.g., plate length, pipe diameter) (m) \\
$\muF$ = Dynamic viscosity of the fluid (Pa$\cdot$s or kg/m$\cdot$s) \\
$\rhoD$ = Fluid density (kg/m$^3$) \\
$\nuF$ = Kinematic viscosity of the fluid ($\nuF = \muF/\rhoD$) (m$^2$/s)

\subsubsection*{Prandtl Number ($\PrNum$)}
\begin{equation}
    \PrNum = \frac{\nuF}{\alphaD} = \frac{\muF \Cp}{\kth}
\end{equation}
where: \\
$\PrNum$ = Prandtl number (dimensionless) \\
$\nuF$ = Kinematic viscosity (m$^2$/s) \\
$\alphaD$ = Thermal diffusivity (m$^2$/s) \\
$\Cp$ = Specific heat capacity (J/kg$\cdot$K) \\
$\kth$ = Thermal conductivity (W/m$\cdot$K)

\subsubsection*{Nusselt Number ($\NuNum$)}
\begin{equation}
    \NuNum_L = \frac{\hC L}{\kth}
\end{equation}
where: \\
$\NuNum_L$ = Nusselt number (dimensionless, based on characteristic length $L$) \\
$\hC$ = Convection heat transfer coefficient (W/m$^2$$\cdot$K) \\
$L$ = Characteristic length (m) \\
$\kth$ = Thermal conductivity of the fluid (W/m$\cdot$K)

\subsubsection*{Stanton Number ($\StNum$)}
\begin{equation}
    \StNum = \frac{\hC}{\rhoD \Vavg \Cp} = \frac{\NuNum}{\ReNum \PrNum}
\end{equation}
where: \\
$\StNum$ = Stanton number (dimensionless)

\subsection{Flat Plate Correlations (Laminar Flow)}
For laminar flow over a flat plate ($ \ReNum_x \lesssim 5 \times 10^5 $):
\subsubsection*{Velocity Boundary Layer Thickness ($\delta$)}
\begin{equation}
    \delta(x) = \frac{5.0 x}{\sqrt{\ReNum_x}}
\end{equation}
where: \\
$\delta(x)$ = Local velocity boundary layer thickness (m) \\
$x$ = Distance from the leading edge (m) \\
$\ReNum_x$ = Local Reynolds number at position $x$ (dimensionless)

\subsubsection*{Local Friction Coefficient ($\Cf,x$)}
\begin{equation}
    \Cf,x = 0.664 \ReNum_x^{-1/2}
\end{equation}

\subsubsection*{Local Nusselt Number ($\NuNum_x$)}
For $\PrNum \ge 0.6$:
\begin{equation}
    \NuNum_x = 0.332 \ReNum_x^{1/2} \PrNum^{1/3}
\end{equation}

\subsubsection*{Thermal Boundary Layer Thickness ($\delta_T$)}
\begin{equation}
    \delta_T \approx \delta \PrNum^{-1/3}
\end{equation}

\subsection{Lumped System Analysis (Transient Convection)}
\subsubsection*{Characteristic Length ($\Lc$)}
\begin{equation}
    \Lc = \frac{V}{\A_s}
\end{equation}
where: \\
$\Lc$ = Characteristic length (m) \\
$V$ = Volume of the body (m$^3$) \\
$\A_s$ = Surface area of the body for heat transfer (m$^2$)

\subsubsection*{Biot Number ($\Bi$)}
\begin{equation}
    \Bi = \frac{\hC \Lc}{\kth_{\text{solid}}}
\end{equation}
where: \\
$\Bi$ = Biot number (dimensionless) \\
$\hC$ = Convection heat transfer coefficient (W/m$^2$$\cdot$K) \\
$\kth_{\text{solid}}$ = Thermal conductivity of the solid body (W/m$\cdot$K)

\subsubsection*{Temperature Distribution with Time}
If $\Bi < 0.1$, the temperature of the body as a function of time can be approximated by:
\begin{equation}
    \frac{\T(t) - \Tinf}{\T_i - \Tinf} = e^{-bt}
\end{equation}
where: \\
$\T(t)$ = Temperature of the body at time $t$ (K or $^\circ$C) \\
$\Tinf$ = Ambient fluid temperature (K or $^\circ$C) \\
$\T_i$ = Initial uniform temperature of the body (K or $^\circ$C) \\
$b = \frac{\hC \A_s}{m \Cp} = \frac{\hC \A_s}{\rhoD V \Cp}$ = Time constant (1/s) \\
$m$ = Mass of the body (kg) \\
$\Cp$ = Specific heat capacity of the body (J/kg$\cdot$K) \\
$\rhoD$ = Density of the body (kg/m$^3$)